% Generated draft: requires mathematical and source review.
\documentclass[11pt]{article}
\usepackage[margin=25mm]{geometry}
\usepackage{fontspec}
\setmainfont{Latin Modern Roman}
\usepackage{amsmath,amssymb,mathtools,unicode-math}
\setmathfont{Latin Modern Math}
\usepackage{enumitem,needspace,xurl,hyperref}
\setlength{\parindent}{0pt}
\setlength{\parskip}{5pt}
\setlength{\emergencystretch}{3em}
\newcommand{\N}{\mathbb{N}}
\newcommand{\Z}{\mathbb{Z}}
\newcommand{\Q}{\mathbb{Q}}
\newcommand{\R}{\mathbb{R}}
\newcommand{\C}{\mathbb{C}}
\newcommand{\F}{\mathbb{F}}
\newcommand{\A}{\mathbb{A}}
\newcommand{\PP}{\mathbb{P}}
\newcommand{\E}{\mathbb{E}}
\newcommand{\eps}{\varepsilon}
\newcommand{\dd}{\,\mathrm{d}}
\begin{document}
\section*{UnsolvedMath ID: 800957}
\textbf{Source problem code:} catalogue entry
\subsection{Definitions and mathematical statement}
Let \(M\) be a closed manifold, meaning compact and without boundary. The supplied excerpt asks about a class of metrics \(\mathcal G\), but does not define that class or its topology. The following expresses the question in the Riemannian interpretation, conditional on specifying \(\mathcal G\), its topology, and sufficient regularity for the geodesic flow.

For each \(g\in\mathcal G\), define its unit tangent bundle by
\[
S_gM=\{(x,v)\in TM:g_x(v,v)=1\}.
\]
A unit-speed closed \(g\)-geodesic is a map \(\gamma:\R/T\Z\to M\), for some \(T>0\), satisfying
\[
\nabla^g_{\dot\gamma}\dot\gamma=0,
\qquad g_{\gamma(t)}(\dot\gamma(t),\dot\gamma(t))=1
\quad\text{for every }t\in\R/T\Z.
\]
Let
\[
P_g=\bigl\{(\gamma(t),\dot\gamma(t)):\ T>0,\ \gamma\text{ is a unit-speed closed }g\text{-geodesic},\ t\in\R/T\Z\bigr\}.
\]
Thus \(P_g\subseteq S_gM\) consists of the unit tangent vectors along all closed geodesics. Define
\[
\mathcal D=\{g\in\mathcal G:\overline{P_g}^{\,S_gM}=S_gM\},
\]
where the closure uses the usual topology of the unit tangent bundle.

\textbf{Question 37.} Is \(\mathcal D\) dense in \(\mathcal G\) for the intended topology on the space of metrics? Equivalently, does every nonempty open subset \(U\subseteq\mathcal G\) contain a metric \(g\) such that every nonempty open subset of \(S_gM\) contains a unit tangent vector to a closed \(g\)-geodesic?

\textit{Unresolved source qualifications.} The excerpt leaves the definition and topology of \(\mathcal G\), the differentiability parameter \(k\), and any standing dimensional or other hypotheses unspecified. Its notation \(SM\) is interpreted here as \(S_gM\) for the metric being considered. Accordingly, the displayed question is a formulation schema pending recovery of those standing assumptions.

The preceding source paragraph describes a question as closely related to the Arnold diffusion conjecture. It also states that, when \(\dim M=2\) and \(k\) is sufficiently large, an answer is negative: it invokes a metric with integrable geodesic flow and certain irrational Liouville tori, and attributes non-transitivity of the geodesic flow of every sufficiently small metric perturbation to KAM theory. The excerpt does not identify the antecedent of that negative answer, define transitivity, specify the perturbation topology or threshold for \(k\), or state the precise hypotheses required for its KAM argument. This contextual assertion is retained as supplied and is not treated as a negative answer to the closed-geodesic density question.
\subsection{Short English statement}
For a closed manifold \(M\), the question asks whether metrics whose closed geodesics have unit tangent vectors dense in the entire unit tangent bundle form a dense subset of the intended metric space \(\mathcal G\). The excerpt omits the definition and topology of \(\mathcal G\) and related standing assumptions, so these must be recovered before the statement is fully specified. The registry labels the problem open; that label is supplied metadata, not an independently verified assessment.
\subsection{Sources}
\begin{itemize}
\item Supplied local UnsolvedMath registry record, permanent ID \(800957\), exact source problem code \texttt{catalogue entry}, titled \emph{Open Problems and Questions about Geodesics --- Question 37}; collection: AMR Open Problem Lists. The registry attributes the record to Burns and Matveev, gives the proposed year as 2019, and labels its status open. These are supplied registry assertions.
\item Burns and Matveev, \emph{Open Problems and Questions about Geodesics} (2019, as identified in the supplied source record), Question 37. Supplied source URL: \url{https://arxiv.org/abs/1308.5417}. The provenance records access on 2026-10-16 and extraction by \texttt{arxiv\_latex}. This is supplied provenance; the source and its version have not been independently consulted or verified here.
\end{itemize}
\end{document}
